\section{Chapter 12}

\begin{verse}
The vision of colors blinds men's eyes\\
The perception of sounds deafens the ears\\
The taste of flavors dulls the mouth\\
Identifying in action darkens the mind\\
Prurient desire destroys the possibility of motion (being tied to the desired object)\\
Therefore: the True Man\\
Does not lose the I in the not-I\\
He keeps out of the exterior, he consists in the interior.
\end{verse}


Aphorisms with possible initiatic significance, in relation to the discipline of avoiding the duress of impressions of the external world (see, also, the ``closing of the gate'') so that to develop a subtle sensibility (that, which in the first lines is said to be muffled or prevented) and to consolidate inner freedom, so that it is equivalent to a gradual movement of the center from the level of the \emph{p'o} soul to that of the \emph{hun} soul. It is indeed Taoist teaching that the natural use of the senses can be detrimental, when it is a question of protecting and preserving intact the ``original simplicity'' (\emph{p'u}): the bursting of the multiplicity of sense appearances into the mind.

As an example of the distance, in the Chinese text, between the literal meaning of the characters and their abstract meaning: literally the penultimate line would sound more or less like this: ``The True Man is for the belly and not for the eye''. If it is that the eye is considered as the gate through which the external world penetrates into the I, it alters it and transports it, muffling the inner senses, while the belly is taken as the ``empty'' part of the body, therefore, transposing (according to the ideas espoused in the preceding chapter), to symbols of the essential and transcendent part of the human being.

The preeminence of the belly over the eye therefore signifies that the True Man does not lose himself in the world that is revealed to the eye, i.e., in the outer world, in the not-I. Something generally unknown in the West, obesity, the ``great belly'' with which ``Immortals'' and Taoist sages, but also certain Chinese and Japanese Buddhas, were often portrayed, has a symbolic character based on the idea just explained: it alludes to the development that has precisely the ``empty'' principle in such beings, or the dimension of ``emptiness'' --- the immaterial or transcendent element. The lower belly --- in Japanese: hara. On its significance of the center in Japanese esoteric schools existing to this day, see K. Durckheim, \emph{Hara, Vital center of man according to Zen} and its problems that are connected there in our essay: The Japanese Hara-theory in its Relation to East and West (in East and West, number 2, 1958).

\paragraph{Chinese text and literal translation}

\textbf{Chapter 12 (\begin{CJK*}{UTF8}{bsmi}第十二章\end{CJK*})}

\columnratio{.4}
\begin{paracol}{2}
\begin{verse}
\begin{CJK*}{UTF8}{bsmi}
五色令人目盲，\\
五音令人耳聾，\\
五味令人口爽，\\
馳騁畋獵令人心發狂，\\
難得之貨令人行妨。\\
是以聖人為腹不為目，\\
故去彼取此。
\end{CJK*}
\end{verse}
\switchcolumn
\begin{verse}
The beautiful colors blind people's eyes,\\
The appealing music stun people's ears\\
The delicious flavors make people's mouth numb,\\
To indulge in hunting makes people's heart wild,\\
To pursue rare treasures makes people's behavior improper.\\
So, as the sage attends to the inner world, not the outer world,\\
Throw away the latter and adopt the former.
\end{verse}
\end{paracol}


\flrightit{Posted on 2020-10-18 by Lao Tzu}

\begin{center}* * *\end{center}

\begin{footnotesize}
\begin{sffamily}
\texttt{Gabe Ruth on 2020-11-18 at 12:41 said:}

Bet my man Chesterton wouldn't have been so down on the wisdom of the East if he'd known about their attitude toward great bellies.


\hfill

\end{sffamily}
\end{footnotesize}

